\section{Introducción}
\label{sec:introduccion}

Un atractor \(\mathcal A\) se denomina \emph{autoexcitado} cuando su cuenca de
atracción intersecta una vecindad de, al menos, un punto de equilibrio. En ese
caso es posible iniciar una trayectoria cerca del equilibrio, descartar el
transitorio y observar cómo la trayectoria se aproxima al atractor. Un
atractor se denomina \emph{oculto} cuando su cuenca permanece separada de
vecindades suficientemente pequeñas de todos los equilibrios
\cite{LeonovKuznetsov2013,KuznetsovEtAl2017}. La palabra ``oculto'' describe,
por tanto, una propiedad de la cuenca: el procedimiento habitual de comenzar
cerca de un equilibrio no proporciona directamente una condición inicial que
alcance el atractor.

La dificultad de localización tiene dos partes. Primero debe encontrarse una
semilla dentro de una región de la cuenca cuya posición se desconoce. Después
debe estudiarse el destino de condiciones iniciales próximas a todos los
equilibrios. La primera parte es un problema de búsqueda; la segunda es un
problema de clasificación de cuencas. Cuando coexisten varios destinos, la
cuenca buscada puede ser estrecha, estar alejada de los equilibrios o presentar
fronteras complejas. Una estrategia de prueba y error depende fuertemente de
las condiciones iniciales elegidas y no indica cómo repetir la búsqueda en
otro conjunto de parámetros.

En un sistema de orden fraccionario aparece una dificultad adicional. Para una
derivada de Caputo, el valor de \(X(t)\) depende de la historia calculada desde
el instante inicial \(t_0\). La condición inicial, el orden \(q\), el paso de
integración, el horizonte temporal y la política de memoria forman parte del
problema dinámico. Reiniciar una continuación utilizando solamente el último
vector de estado cambia la historia y, por ello, cambia el problema de valor
inicial. El orden \(q\) también modifica la estabilidad local: la condición
usual \(\Re\lambda<0\) se sustituye por el criterio angular de Matignon
\cite{Matignon1996}.

Este reporte desarrolla dos rutas complementarias:

\begin{enumerate}
  \item para el Chua no suave, la semilla se obtiene mediante una función
  descriptiva sesgada y se transporta con una continuación causal que conserva
  una sola historia de Caputo;
  \item para el Chua con arctangente, los parámetros se exploran primero en
  orden entero y la condición inicial se refina después mediante problemas de
  valor inicial de Caputo independientes.
\end{enumerate}

En ambos casos, los sondeos de cuenca comienzan en el mismo instante de memoria.
El caso no suave conserva la regla de completar el primer radio con contactos.
En \texttt{c590} se añadió una auditoría macroscópica en \(r=2\) después del
primer contacto observado en \(r=1\). El mayor radio local sin contactos y los
bloques macroscópicos se reportan por separado. El alcance \emph{local} se fija
por el contrato predeclarado. Los cálculos cubren radios, geometrías,
cantidades de muestras y horizontes declarados. Una caracterización global
exigiría describir la cuenca completa en un dominio no acotado y con resolución
arbitrariamente fina, tarea que es especialmente costosa cuando cada
trayectoria de Caputo acumula toda su historia.

\section{Antecedentes y métodos de referencia}
\label{sec:antecedentes}

\subsection{Leonov y Kuznetsov: método en orden entero}

Leonov, Kuznetsov y colaboradores formularon un procedimiento
analítico--numérico para localizar atractores ocultos en sistemas tipo Lur'e
\cite{LeonovKuznetsov2013,KuznetsovEtAl2017}. El procedimiento entero puede
resumirse en cuatro operaciones:

\begin{enumerate}
  \item separar el sistema en un bloque lineal y una no linealidad escalar;
  \item aproximar la respuesta de la no linealidad por su componente
  fundamental y resolver una condición de balance armónico;
  \item construir una condición inicial a partir del modo oscilatorio del
  sistema lineal auxiliar;
  \item transportar esa condición inicial, mediante continuación, hasta el
  sistema no lineal original.
\end{enumerate}

La función descriptiva determina una frecuencia, una ganancia equivalente y
una amplitud. Esos valores se utilizan para construir la semilla. La
clasificación del atractor se obtiene posteriormente al estudiar la cuenca
alrededor de los equilibrios. Esta separación entre \emph{localización} y
\emph{clasificación de cuenca} se conserva en toda la metodología propuesta.

\subsection{Danca: Chua fraccionario no suave}

Danca estudió atractores ocultos en varios sistemas fraccionarios, incluido un
sistema de Chua con saturación \cite{Danca2017}. El artículo usa derivadas
de Caputo de orden conmensurable, clasifica los equilibrios mediante Matignon e integra
el problema con un predictor--corrector Adams--Bashforth--Moulton. Para el
ejemplo de Chua se reportan
\[
q=0.9998,\qquad
(\alpha,\beta,\gamma)=(8.4562,12.0732,0.0052),\qquad
(m_0,m_1)=(-0.1768,-1.1468).
\]
La verificación descrita en la referencia usa 200 trayectorias alrededor de un
equilibrio externo con radio \(\delta=0.01\) y emplea la simetría para el
equilibrio opuesto.

La reconstrucción realizada en este trabajo conserva la ecuación, los
parámetros, el orden y el integrador de la referencia. El artículo no
proporciona la condición inicial exacta de la figura del atractor ni los
valores intermedios de la construcción armónica
\((\omega_0,k,A_0)\). Por ello, la reconstrucción utiliza la semilla obtenida
al aplicar el cierre armónico entero a los parámetros publicados y declara
esta elección de manera explícita. La ruta propuesta usa los mismos parámetros
nominales, pero calcula una nueva semilla sesgada y aplica una continuación
compatible con la memoria de Caputo.

\subsection{Wu et al.: Chua fraccionario con arctangente}

Wu et al. propusieron un sistema de Chua fraccionario con no linealidad
arctangente \cite{Wu2023ArctanChua}. La referencia emplea una función
descriptiva, una recurrencia local de descomposición de Adomian y las
condiciones iniciales opuestas
\[
X_{0,+}=(13.8,\;0.7093,\;-19.8768)^{\T},
\qquad X_{0,-}=-X_{0,+}.
\]
Los parámetros publicados se expresan mediante
\[
h(x)=mx+(n-m)\arctan(x),\qquad
m=0.4,\qquad n=-1.1585,
\]
de modo que, en la notación empleada aquí,
\[
a_1=0.4,\qquad a_2=n-m=-1.5585,\qquad \rho=1,\qquad q=0.99.
\]

La reconstrucción bibliográfica reproduce la recurrencia ADM local y las dos
condiciones iniciales publicadas. La ruta propuesta para el sistema con
arctangente es independiente: explora un dominio acotado de parámetros en
\(q=1\), selecciona un caso mediante reglas reproducibles y refina la semilla
con problemas de Caputo de memoria completa.

\subsection{Relación entre las tres referencias}

Las referencias cumplen funciones diferentes dentro del estudio:

\begin{table}[H]
\centering
\caption{Papel de los trabajos de referencia en la metodología.}
\label{tab:papel_referencias}
\begin{tabularx}{\linewidth}{L{0.22\linewidth}Y Y}
\toprule
Referencia & Elemento recuperado & Uso en este reporte\\
\midrule
Leonov--Kuznetsov
& Función descriptiva, construcción modal y continuación en \(q=1\)
& Caso límite y fundamento de la ruta de localización\\
Danca
& Chua no suave, Caputo, Matignon y ABM
& Reconstrucción bibliográfica y sistema base de la ruta sesgada\\
Wu et al.
& Chua con arctangente, condiciones iniciales y ADM local
& Reconstrucción bibliográfica y comparación con la ruta independiente
\texttt{c590}\\
\bottomrule
\end{tabularx}
\end{table}

La jerarquía bibliográfica del resto del reporte es deliberada. La definición,
la localización y los parámetros de cada atractor se atribuyen a artículos de
investigación
\cite{LeonovKuznetsovVagaitsev2011,Leonov2013HiddenAttractors,Danca2017,KuznetsovEtAl2017,Wu2023ArctanChua};
los métodos geométricos se remiten también a sus trabajos originales
\cite{DudkowskiPrasadKapitaniak2017,GuanXie2021Connecting,GodaraEtAl2018CriticalSurfaces}.
Los libros compilados
\cite{PhamEtAl2018HiddenAttractors,PhamVolosKapitaniak2017Systems,WangKuznetsovChen2021Multistability}
se usan para ordenar familias y encontrar antecedentes. Ningún parámetro,
condición inicial o reclamo de ocultedad se sostiene únicamente en un capítulo
recopilatorio o en un resumen automático.

El orden lógico del reporte sigue esa separación: primero se definen los
sistemas y los fundamentos; después se derivan las ecuaciones comunes; luego
se describen los protocolos bibliográficos y las dos rutas propuestas; por
último se comparan los resultados obtenidos.
